\documentclass[10pt]{article}
\usepackage[margin=1in]{geometry}
\usepackage[T1]{fontenc}
\usepackage[utf8]{inputenc}
\usepackage{amsmath,amssymb}
\usepackage{booktabs}
\usepackage{tabularx}
\usepackage{enumitem}
\usepackage{microtype}
\setlength{\parindent}{15pt}
\setlength{\parskip}{0pt}
\setlist[itemize]{leftmargin=2em, itemsep=0.55em, topsep=0.4em}

\title{Free-Space Magnetic Levitation\\A Two-Needle Hand Knitting Mechanics\\Framework}
\author{}
\date{}

\begin{document}
\maketitle

\begin{abstract}
Existing magnetically driven knitting systems constrain needles mechanically within a machine bed, with magnetic actuation controlling only vertical displacement. This paper proposes an exploratory, un-prototyped architecture: a two-needle system modeled on hand knitting mechanics, where full six-degree-of-freedom (6-DOF) needle actuation is hypothesized purely through magnetic means without physical guides. The needle design utilizes a conductive core for eddy current torque generation paired with a non-continuous magnetic strip for discrete pole levitation control. Furthermore, the framework introduces textile physics into the control laws, tracking variations in mass, moment of inertia, and gyroscopic cross-coupling induced by the changing orientation angle $\theta$ of the dynamic V-shaped needle configuration and fluctuating loop counts ($n$). An active, bidirectional motorized yarn dispenser is proposed as a mechanism to manage loop tension and eliminate slack throughout intermediate stitch cycles.
\end{abstract}

\section{Introduction}
Automated knitting has long been dominated by industrial machines that use mechanical cam systems to drive arrays of needles. These systems are effective but mechanically complex, noisy, and prone to needle breakage due to physical contact between components.

Recent work has explored replacing mechanical cam drives with magnetic actuation. Zhang et al. [1] proposed a non-contact magnetic array needle driving method where permanent magnet needles are driven vertically by an electromagnet array. Subsequent work [2, 3, 4] developed this approach further, optimizing needle gauge, magnetic shielding between needles, and control strategies including PID and sliding mode controllers. However, all of these systems share a fundamental architectural constraint: needles are mechanically guided within a machine bed, and magnetic actuation controls only the vertical (axial) displacement of each needle. The architecture remains that of an industrial weft knitting machine with a magnetic upgrade.

This paper proposes a theoretical departure from that architecture. Rather than scaling up an industrial machine, the goal is to evaluate the conceptual mechanics of hand knitting --- two needles, free-space actuation, full 6-DOF control --- using a solenoid array and purpose-designed needles. C The main contributions of this exploratory work are:
\begin{itemize}
  \item A two-needle free-space magnetic actuation architecture modeled on the dynamic, non-parallel paths of hand knitting.
  \item A purpose-designed needle where geometric and material configurations serve specific degrees of freedom, now augmented by variables tracking distributed yarn mass.
  \item A mechanical tracking model incorporating the changing intersection angle ($\theta$) of the needles to counter gyroscopic effects, paired with an active, bidirectional motorized yarn dispenser for real-time tension regulation.
\end{itemize}

\section{System Overview}
The proposed system consists of two identical actuator units, each comprising a dense solenoid array and a knitting needle. The two units are positioned facing each other, separated by an insulating boundary, in a configuration that mirrors the two-handed mechanics of hand knitting.

Each solenoid array consists of a $p\times q$ grid of individually addressable miniature solenoids with silicon-steel laminated cores. Each knitting needle sits above its respective array and is modeled to be actuated purely through magnetic forces and torques generated by the array. No mechanical guides constrain the needle's motion. The fabric is initially cast onto the needle shafts manually or via an auxiliary startup sequence, establishing a pre-loaded baseline state before magnetic cycling begins.

The knitted fabric is collected in a wooden tray positioned below the work area. The tray height is set according to the estimated breaking stress of the needle and the linear mass density of the knitting pattern, ensuring the needle is not overloaded by the weight of accumulated fabric.

\section{Needle Design}
The needle is designed so that each component directly serves one or more degrees of freedom. No element is arbitrary.

\subsection{Geometry}
The needle has a cylindrical shaft. The leading end is conical, optimized for yarn manipulation during loop formation. The trailing end is flat, which maximizes the surface area available for eddy current generation and therefore torque transfer.

\subsection{Aluminum Core}
The core of the needle is aluminum. Aluminum is non-ferromagnetic, meaning it does not retain or redirect magnetic fields passively. In regions of the needle where no magnetic strip is present, the aluminum core is magnetically inert to static or slowly varying fields. This eliminates the need for additional passive shielding on the shaft.

Under a time-varying magnetic field, eddy currents are induced in the aluminum core. These eddy currents interact with the driving magnetic field to produce torque, which controls the rotational degrees of freedom. The flat trailing end maximizes this effect by increasing the effective cross-sectional area for eddy current circulation.

\subsection{Non-Continuous Magnetic Strip}
A magnetic strip runs along the shaft of the needle, interrupted periodically (e.g., every inch for an inch). This non-continuous configuration creates discrete magnetic poles along the needle's length. These poles are the primary interaction points with the solenoid array for levitation and translational control, allowing spatial control of the magnetic interaction and smoother transitions between solenoid activations.

\subsection{No Mu-Metal Coating}
Mu-metal was considered for passive shielding of the non-strip regions of the shaft. It was rejected on three grounds: cost, mechanical durability under repetitive yarn friction, and physical redundancy. Since aluminum is non-ferromagnetic, the non-strip shaft regions are already magnetically inert without additional coating. Cross-interference from these regions is managed through PWM frequency control and solenoid sequencing rather than passive shielding coatings.

\section{Degrees of Freedom and Control Architecture}
The needle requires full 6-DOF actuation. Three translational degrees of freedom and three rotational degrees of freedom are identified, each mapped to a specific control mechanism.

\subsection{Coordinate System}
A right-handed Cartesian coordinate system is defined with its origin at the center of the solenoid array surface. The Y-axis is normal to the array surface (vertical). The X and Z axes span the horizontal plane of the array.

\subsection{DOF-to-Mechanism Mapping}
\begin{table}[htbp]
\centering
\begin{tabularx}{\linewidth}{l X X}
\toprule
DOF & Description & Control Mechanism \\
\midrule
X & Horizontal position (left-right) & Selective solenoid activation \\
Z & Horizontal position (forward-back) & Selective solenoid activation \\
Y & Vertical height (levitation) & PWM-controlled lifting force \\
XY plane & Tilt about Z-axis ($\theta_{xy}$) & Eddy current torque \\
YZ plane & Tilt about X-axis ($\theta_{yz}$) & Eddy current torque \\
ZX plane & Rotation about Y-axis ($\theta_{zx}$) & Eddy current torque \\
\bottomrule
\end{tabularx}
\caption{Degree of freedom mapping for needle actuation.}
\end{table}

\subsection{Translational Control (X, Z)}
Horizontal needle positioning is achieved by selectively activating solenoids ahead of and behind the needle's current position, analogous to the operating principle of a linear motor. Real-world hand knitting trajectories are highly structured, cyclical, and repetitive. Thus, instead of an open-ended, general-purpose 6-DOF path planner, the system operates on pre-defined parametric trajectories, significantly constraining the search space and computing overhead.

\subsection{Vertical Control (Y) and Earnshaw's Theorem}
Earnshaw's theorem states that a static magnetic field alone cannot stably levitate a passive magnetic object in free space. The system addresses this through dynamic PWM control rather than static field equilibrium. Because the solenoid currents are continuously varied, the field is never static.

During active knitting, the vertical levitation condition requires that the vertical component of the magnetic lifting force $F_y(t)$ exceed the combined gravitational load and the downward mechanical forces exerted by the yarn loops:
\begin{equation}
F_y(t)>m_{\mathrm{eff}}(n)\cdot g+F_{\mathrm{yarn},y}(d,\theta)
\tag{1}
\end{equation}
where $m_{\mathrm{eff}}(n)$ is the variable effective mass of the needle carrying $n$ active loops of yarn, and $F_{\mathrm{yarn},y}$ is the vertical component of the yarn tension vector, which is a function of the needle separation distance $d$ and the relative intersection angle $\theta$.

\section{Textile Dynamics and Mathematical Framework}
Knitting mechanics dictate that the needles are rarely parallel; instead, they intersect at a continuously changing angle $\theta(t)$, forming a dynamic V-shaped configuration. Furthermore, the yarn loops act as a distributed mass that constantly transfers between the two needles.

\subsection{Dynamic Mass and Moment of Inertia Model}
Let $\lambda$ represent the linear mass density of the yarn (e.g., in g/m or tex), and $t_y$ represent the physical thickness (diameter) of the yarn strand. The total number of active loops across both needles is assumed to be constant during a stable knitting sequence, denoted by $N_{\mathrm{total}}$. At any given timestamp, let $n_1(t)$ be the number of loops on Needle 1, and $n_2(t)$ be the number of loops on Needle 2, such that:
\begin{equation}
n_1(t)+n_2(t)=N_{\mathrm{total}}
\tag{2}
\end{equation}
If $L_{\mathrm{loop}}$ is the average length of yarn wrapped around the needle shaft to form a single loop, the mass contribution of a single loop is $m_{\mathrm{loop}}\approx\lambda\cdot L_{\mathrm{loop}}$. The effective mass of a given needle carrying $n$ loops is modeled as:
\begin{equation}
m_{\mathrm{eff}}(n)=m_{\mathrm{needle}}+n(\lambda\cdot L_{\mathrm{loop}})
\tag{3}
\end{equation}
Because these loops are distributed along the length of the shaft occupying a zone of thickness $n\cdot t_y$, they modify the mass moment of inertia tensor $I_{\mathrm{eff}}(n,\theta)$ relative to the center of mass. Consequently, the eddy current torque ($\tau_{\mathrm{eddy}}$) required to tilt or orient the needles must adjust dynamically as loops transfer ($n$ shifts) and as the angle $\theta$ changes.

\subsection{Gyroscopic Effect Compensation}
When the needles are tilted to alter the V-configuration angle $\theta(t)$ at an angular velocity $\omega_{\mathrm{precession}}$, the combination of the needle's longitudinal rotation and its shifting orientation introduces an unexpected gyroscopic torque:
\begin{equation}
\tau_{\mathrm{gyro}}=\omega_{\mathrm{precession}}\times\left(I_{\mathrm{eff}}(n,\theta)\cdot\omega\right)
\tag{4}
\end{equation}
Because the needle floats freely in space, this cross-coupling torque could easily induce a destabilizing wobble along its longitudinal axis. To preserve stable levitation, the PWM-driven magnetic field must actively apply an equal and opposite corrective torque. The net required control torque vector $\tau_{\mathrm{required}}(t)$ is formulated as:
\begin{equation}
\tau_{\mathrm{required}}(t)=I_{\mathrm{eff}}(n,\theta)\frac{d\omega}{dt}+\omega\times\left(I_{\mathrm{eff}}(n,\theta)\cdot\omega\right)
\tag{5}
\end{equation}
where the second term represents the predictive gyroscopic compensation factor calculated from the pre-defined trajectory paths.

\subsection{Intermediate Stitch State and Active Yarn Dispenser Control}
Between complete stitches, the loops rest uniformly on the needle shafts. However, during an active intermediate stitch phase, the needles separate to pull a new loop through. During this bridging phase, the yarn undergoes high tensile strain and acts as a mechanical tether connecting the two needles. The required yarn length $L_{\mathrm{required}}$ across this dynamic V-shaped gap is a trigonometric function of the instantaneous distance $d$ between the needle tips, the angle $\theta$, and the existing loop geometry:
\[
L_{\mathrm{required}}(t)=d\left[\csc\left(\frac{\theta}{2}\right)+1\right]+L_{\mathrm{loop}}+x
\]
where $x$ is a structural allowance factor accounting for tip geometry and fabric tortuosity.

Yarn possesses high tensile strength under tension (pull) but zero compressive stiffness (push). If $d(t)$ or $\theta(t)$ decreases rapidly, the yarn will immediately slacken, risking snags or destabilizing the needle's magnetic equilibrium. To shift the responsibility of tension control away from the un-guided yarn filament, an active, bidirectional motorized yarn dispenser is integrated as a synchronized slave axis.

Because the master controller executes pre-defined parametric trajectories, the instantaneous values of $d(t)$, $\theta(t)$, and their time-derivatives $\dot{d}(t)$, $\dot{\theta}(t)$ are known a priori. The yarn dispenser motor torque ($\tau_{\mathrm{dispenser}}$) and velocity are governed by the rate of change of the required length:
\begin{equation}
\frac{dL_{\mathrm{required}}}{dt}=\dot{d}\,\csc\left(\frac{\theta}{2}\right)-d\dot{\theta}\,\csc\left(\frac{\theta}{2}\right)\cot\left(\frac{\theta}{2}\right)
\tag{7}
\end{equation}
When $\dot{L}_{\mathrm{required}}>0$ (Pull Phase), the dispenser rotates clockwise to feed out yarn at a controlled tension. When $\dot{L}_{\mathrm{required}}<0$ (Push/Retract Phase), the dispenser actively reverses direction (rotating counter-clockwise), reeling the excess yarn back in at the exact matching rate to prevent any structural slack without requiring real-time computer vision tracking loops.

\section{Yarn Interaction and AI Monitoring}
Yarn properties --- tension, friction, and elasticity --- vary dynamically during knitting and directly affect loop formation quality. Rather than using contact-based tension sensors, this system employs a camera-based AI monitoring approach. A vision system observes the yarn and needle positions in real time and adjusts the solenoid array PWM parameters accordingly.

This approach avoids the complexity of integrating physical sensors into the work area while providing holistic feedback that captures multiple yarn variables simultaneously. The AI system is also responsible for active Y-axis stabilization, compensating for Earnshaw instability through continuous field adjustment based on observed needle position.

\section{Challenges and Mitigations}
\subsection{Earnshaw Instability}
Static magnetic levitation of a passive object is inherently unstable. This is mitigated through dynamic PWM actuation and AI-based visual feedback for active stabilization.

\subsection{Magnetic Cross-Talk}
Adjacent solenoids can interfere with each other's fields. Mitigation strategies include: PWM sequencing to prevent conflicting adjacent activations, array layout symmetry, and natural field attenuation with distance. The non-ferromagnetic aluminum shaft provides passive inertness in non-strip regions without requiring mu-metal coating.

\subsection{Thermal Management}
Continuous solenoid operation generates heat. Mitigation strategies include: PWM duty cycle optimization (solenoids cool during inactive periods), silicon-steel laminated cores (reduced core losses), and increased spacing between adjacent solenoids. The external power supply prevents heat transfer into the work area.

\subsection{Projectile Effect}
If a ferromagnetic object enters the active magnetic field, it will experience a strong attractive force and may be launched unpredictably. The system must be operated in a controlled environment free of ferromagnetic materials.

\subsection{Strong Magnetic Field}
Prolonged human exposure to the operating magnetic field should be avoided. Appropriate shielding of the external environment is recommended.

\section{Questions and Answers}
Won't gravity restrict the needle from levitating?

The gravitational force on the needle does not prevent levitation. The vertical component of the magnetic lifting force is calibrated to satisfy $F_y(t)>m_{\mathrm{eff}}(n)g+F_{\mathrm{yarn},y}$, directly counteracting the needle's total mass and yarn tension vectors.

What happens to the knitted fabric?

The knitted fabric is collected in a wooden tray positioned below the work area. The tray height is determined by the breaking stress of the needle and the linear mass density $\lambda$ of the chosen knitting pattern. If $W$ is the maximum weight the needle can support without bending, the tray is placed at a conservative distance below the needle's lowest operating position to limit the suspended fabric weight.

\section{Limitations}
\begin{itemize}
  \item The mathematical framework presented here is strictly theoretical. No physical prototype has been constructed, and the coupled equations for electromagnetic levitation, gyroscopic torque cancellation.
  \item The efficiency of the thin-film Seebeck matrix and the activation thresholds of the non-mechanical sorption refrigeration channels are highly speculative; miniature solenoids may not generate high enough localized temperature gradients to run these cycles effectively without stalling.
  \item The predictive active yarn dispenser model assumes perfect adherence to the pre-defined parametric trajectories. Any unmodeled slip, yarn friction variation, or external disturbance could introduce slack or over-tension, as real-time camera-based AI vision tracking has been paused in this iteration of the architecture.
\end{itemize}

\section{Future Work}
\begin{itemize}
  \item \textbf{Prototype Development:} Construction and empirical validation of a single-needle scaled-down solenoid array to measure actual $I^2R$ thermal profiles.
  \item \textbf{Numerical Simulation:} Finite Element Analysis (FEA) tracking the thermal-resistance coupling of the copper coils and the simultaneous multi-axis eddy current torques.
  \item \textbf{Textile Characterization:} Experimental testing of yarn friction against an aluminum needle tip to calibrate the structural allowance variable $x$ in the intermediate stitch equation.
\end{itemize}

\section*{References}
\begin{enumerate}[label={[\arabic*]}, leftmargin=2em, itemsep=0.5em]
  \item Zhang, C., Li, D., Zuo, X., \& Zhou, X. (2019). Magnetic array knitting needle driving system design. \textit{IOP Conference Series: Materials Science and Engineering}, 612(3), 032026.
  \item Zuo, X., Zhang, C., Xiong, T., Yin, W., Li, M., Zhang, C., Wu, X., \& Zhu, L. (2021). Maglev weft knitting needle driving modeling and PID controller design. \textit{The Journal of The Textile Institute}, 113(8), 1715--1722.
  \item Xiong, T., Peng, Y., Zuo, X., Zhang, C., Zhang, C., Zhang, L., \& Li, H. (2023). Magnetic field analysis and optimization of the gauge of hybrid maglev needles. \textit{Applied Sciences}, 13(3), 1257.
  \item Zhang, D., Sheng, X., Gao, P., Wang, C., \& Shi, Y. (2025). Structure design and control of magnetic levitated needle driving system. In \textit{Intelligent Robotics and Applications, ICIRA 2024}, Lecture Notes in Computer Science, vol. 15207. Springer, Singapore.
\end{enumerate}

\end{document}
